% GPT-6 Astra Pro; effort not exposed; /root/replacement; 22 November 2026.
\section{HD-F1-R10: an extreme pure cohomology sheaf across a gap gives a trace kernel}
\label{sec:followup-rm-gap-obstruction}

Let $E$ be any actual perfect object on the genuine real-quadratic
fourfold $X$ of R07, with the full balanced character
\[
 \operatorname{ch}(E)=N e^{-B}v_q,\qquad
 N\ne0,\quad q>0,\quad B\in\operatorname{NS}(X)_{\mathbf Q}.
 \tag{F1R10.1}
\]
Write $a$ and $b$ for its least and greatest nonzero standard
cohomology degrees. Each of the following actual cohomology
patterns forces an invariant full-trace kernel of dimension at least
two in degree two for every coherent faithful finite point action and every
lift of $E$:
\[
 \begin{array}{ll}
 H^{a+1}(E)=0,&
 H^a(E)\text{ has a nonzero pure-character orbit of R09};\\[2pt]
 H^{b-1}(E)=0,&
 H^b(E)\text{ has a nonzero pure-character orbit of R09}.
 \end{array}
 \tag{F1R10.2}
\]
The phrase in the second column means exactly (F1R9.2), or the
character obtained from it by the finite compositions of actual
functors specified in R08--R09. Thus it includes an extreme line
bundle, a simple semihomogeneous bundle, or a skyscraper sheaf,
with its indicated cohomological shift. No preservation of these
cohomology sheaves by a chosen linearization is required.
The other cohomology sheaves can be arbitrary.

In particular every actual two-stage object
\[
 H^0(E)=A,\quad H^m(E)=T,\quad
 H^j(E)=0\ (j\ne0,m),\qquad m\ge2,
 \tag{F1R10.3}
\]
with at least one of the two coherent sheaves having the indicated
pure-character orbit, fails weak $G$-semiregularity for every
such lift when its whole character is (F1R10.1). This includes
all actual nonzero extension classes in
$\operatorname{Ext}^{m+1}_X(T,A)$, not only the split triangle.
It therefore excludes the pure-endpoint part of the
amplitude $[0,2]$ cone construction in DEG\_L S-27 by an actual
map of self-Ext groups. It supplies no claim about the
remaining cones whose two extreme characters both fail the
pure-character condition.

\paragraph{1. The canonical truncation triangles and exact injections.}
First consider the least cohomology degree in (F1R10.2). The
canonical standard truncation triangle is
\[
 Q:=H^a(E)[-a]\xrightarrow{i}E\longrightarrow
 R:=\tau^{\ge a+2}E\longrightarrow Q[1].
 \tag{F1R10.4}
\]
The assumed zero adjacent cohomology makes this the actual
truncation triangle. Since $Q\in D^{\le a}$ and
$R[1]\in D^{\ge a+1}$, standard t-structure orthogonality gives
$\operatorname{Hom}(Q,R[1])=0$. The long exact Hom sequence
therefore supplies the injection
\[
 \operatorname{Ext}^2_X(Q,Q)\lhook\joinrel\longrightarrow
         \operatorname{Hom}_X(Q,E[2]),\qquad
 u\longmapsto i[2]u.
 \tag{F1R10.5}
\]
For the greatest degree the canonical triangle is instead
\[
 A:=\tau^{\le b-2}E\longrightarrow E
       \xrightarrow{p}Q:=H^b(E)[-b]\longrightarrow A[1].
 \tag{F1R10.6}
\]
Now $A\in D^{\le b-2}$ and $Q[1]\in D^{\ge b-1}$, so
$\operatorname{Hom}(A,Q[1])=0$. Applying
$\operatorname{Hom}(-,Q[2])$ gives the second exact injection
\[
 \operatorname{Ext}^2_X(Q,Q)\lhook\joinrel\longrightarrow
         \operatorname{Hom}_X(E,Q[2]),\qquad
 u\longmapsto up.
 \tag{F1R10.7}
\]
These vanishing statements follow from the actual cohomology
gap; they are not additional unsolved Ext conditions.

\paragraph{2. A nonzero invariant evaluation on the coupled object.}
We first strengthen the quotient calculation of R09 quantitatively.
For any surjection $\pi:U_1\oplus U_2\to W$, $\dim W=4$, put
$r_i=\operatorname{rank}(\pi|_{U_i})$. The exact ranks on
$K_i=\ker\omega_i$ are
\[
 \operatorname{rank}\bigl((\bigwedge^2\pi)|_{K_i}\bigr)
 =\begin{cases}
 0,&r_i=0,1,\\
 1,&r_i=2,\\
 3,&r_i=3,\\
 5,&r_i=4.
 \end{cases}
 \tag{F1R10.8a}
\]
For rank two, the determinant form is decomposable and cannot
be proportional to the nondegenerate $\omega_i$, so its
restriction to $K_i$ is nonzero. For rank three, the kernel
of the exterior square is $(\ker\pi_i)\wedge U_i$, of dimension
three. Nondegeneracy of $\omega_i$ makes its restriction to
this kernel nonzero, so its intersection with $K_i$ has
dimension two. This gives rank three. The other entries
are immediate.

Surjectivity gives $r_1+r_2\ge4$. If either rank is at least
three, the image of $K$ already has dimension at least three.
Otherwise both ranks are two, their images are complementary
two-planes in $W$, and their two exterior-volume lines are
independent. Consequently
\[
       \operatorname{rank}\bigl((\bigwedge^2\pi)|_K\bigr)\ge2.
 \tag{F1R10.8b}
\]
The exact shears and transports of R09 carry this to the same
lower bound for
$c^2_{\operatorname{ch}(Q)}|_{\ker c^2_{\operatorname{ch}(E)}}$.
In particular the direct-summand trace-kernel bound of R09
also improves to two.

The quotient lemma and exact character transports of R09 apply
to the indicated $Q$. They produce a particular class
\[
 z\in HH^2(X),\qquad
 c^2_{\operatorname{ch}(E)}(z)=0,\qquad
 c^2_{\operatorname{ch}(Q)}(z)\ne0.
 \tag{F1R10.8}
\]
Consequently $\operatorname{ev}_Q(z)\ne0$ by the full
trace--Atiyah identity. In the first triangle, naturality
identifies its image under (F1R10.5) with
\[
 i[2]\operatorname{ev}_Q(z)
               =\operatorname{ev}_E(z)i.
 \tag{F1R10.9}
\]
In the second triangle, its image under (F1R10.7) is
\[
 \operatorname{ev}_Q(z)p
               =p[2]\operatorname{ev}_E(z).
 \tag{F1R10.10}
\]
The appropriate injection proves that this composite is
nonzero. Thus $\operatorname{ev}_E(z)\ne0$, while its
full trace is zero by (F1R10.8). The characteristic action
commutes with every specified point linearization, so this
is a nonzero element of
$\operatorname{Ext}^2_X(E,E)^G$. The exact identity-sector
forgetful isomorphism already proved in R05 identifies its
kernel with the trace kernel of the equivariant lift.
For the dimension assertion, apply the same maps to the entire
space $\ker c^2_{\operatorname{ch}(E)}$. Its evaluation on $Q$
has rank at least two by (F1R10.8b) and the trace--Atiyah
identity. The injection (F1R10.5) or (F1R10.7) therefore forces
rank at least two for its evaluation on $E$. All these values
are invariant and have zero full trace. This proves the
dimension bound for the original weak trace map.

\paragraph{3. The actual two-sheaf cone and the framework maps.}
For (F1R10.3) the canonical triangle is
$A\to E\to T[-m]\to A[1]$. Both injections above use the
single explicit vanishing
\[
 \operatorname{Hom}(A,T[-m][1])
                    =\operatorname{Ext}^{\,1-m}_X(A,T)=0,
 \tag{F1R10.11}
\]
which holds for coherent sheaves and $m\ge2$. Thus the
argument treats every actual coupling morphism
$T[-m]\to A[1]$, including nonzero ones. In the
Grothendieck group its whole character is the realization
of the actual triangle relation
$[E]=[A]+(-1)^m[T]$; GRR and the Adams/point-biextension
maps are exactly (F1R9.9)--(F1R9.12). The Ext conclusion
uses the two canonical truncation arrows and the displayed
long exact sequences, followed by the same published
Buchweitz--Flenner characteristic map and Perry
identity-sector comparison. These maps give the stated
extension beyond actual direct summands without asserting
kernel persistence for arbitrary cones.

\paragraph{4. The adjacent two-cohomology case has an exact cancellation bound.}
There is also a quantitative restriction on every actual object
with $H^0(E)=A$, $H^1(E)=T$ and no other cohomology, whose whole
character is (F1R10.1), and with at least one endpoint having
the pure-character orbit of R09. Write its actual truncation
triangle and connecting class as
\[
 A\xrightarrow{i}E\xrightarrow{p}T[-1]
       \longrightarrow A[1],\qquad
 \delta\in\operatorname{Ext}^2_X(T,A),\qquad
 h=\dim\operatorname{Hom}_X(A,T).
 \tag{F1R10.12}
\]
For every coherent faithful finite point action and lift of $E$,
\[
        \dim\ker\sigma^2_{\widetilde E}\ge\max(0,2-h).
 \tag{F1R10.13}
\]
Thus every such actual adjacent object with $h\le1$ is excluded;
weak semiregularity requires $h\ge2$. This is a necessary
restriction and excludes no object with $h\ge2$ solely from
that dimension.

The precise kernels replacing (F1R10.5) and (F1R10.7) are
the images of the actual Yoneda maps
\[
 \begin{aligned}
 Y_A:\operatorname{Hom}(A,T)&\longrightarrow
                     \operatorname{Ext}^2(A,A),
             &f&\longmapsto\delta f,\\
 Y_T:\operatorname{Hom}(A,T)&\longrightarrow
                     \operatorname{Ext}^2(T,T),
             &f&\longmapsto f[2]\delta.
 \end{aligned}
 \tag{F1R10.14}
\]
These follow from the same two long exact Hom sequences;
$T[-1]\to A[1]$ corresponds after shifting to
$\delta:T\to A[2]$. Each Yoneda image has dimension at most $h$.
The evaluation of $K_E=\ker c^2_{\operatorname{ch}(E)}$ on the
pure endpoint has image of dimension at least two by
(F1R10.8b). Passing through $i[2]$ or precomposing with $p$
can lose at most the $h$ dimensions in (F1R10.14).
Naturality (F1R10.9) or (F1R10.10) gives an image of
dimension at least $2-h$ in the invariant trace kernel of $E$.
The identity-sector comparison proves (F1R10.13).

The same maps give a nonnumerical necessary condition: every
weakly semiregular lift satisfies
\[
 \operatorname{ev}_A(K_E)\subseteq\operatorname{im}Y_A,\qquad
 \operatorname{ev}_T(K_E)\subseteq\operatorname{im}Y_T.
 \tag{F1R10.15}
\]
Indeed invariant evaluation on $E$ vanishes on $K_E$ under
an injective weak trace. This does not assert a common
$f$ for the two endpoint evaluations; the two inclusions
are proved for the actual connecting morphism $\delta$.
